




\smallskip
\noindent\textbf{Traditional approach.}
Dynamic methods infer candidate invariants from traces: Daikon~\cite{Daikon} filters templates against observed tests, Sharma et al.~\cite{A_Data_Driven_Approach_for_Algebraic_Loop} add SMT validation, and DIG~\cite{dig} infers equations and polyhedra directly. Their results depend on test coverage. Static methods avoid this dependence through constraint solving~\cite{DBLP:conf/cav/ColonSS03,DBLP:journals/pacmpl/LiuFYSL22}, abstract interpretation~\cite{Abstract_Interpretation_A_Unified_Lattice_Model_for_Static_Analysis_of_Programs_by_Construction_or_Approximation_of_Fixpoints,Automatic_Discovery_of_Linear_Restraints_among_Variables_of_a_Program}, interpolation~\cite{FiB_Squeezing_loop,DBLP:conf/cade/GanDXZKC16,Nonlinear_Craig_Interpolant_Generation}, recurrence analysis~\cite{ISSACDeepakKapur,DBLP:journals/jsc/Rodriguez-CarbonellK07,DBLP:journals/pacmpl/CyphertK24}, and abductive or shape reasoning~\cite{Inductive_Invariant_Generation_via_Abductive_Inference,Backward_Symbolic_Execution_with_Loop_Folding,SymInfer_Inferring_Numerical_Invariants_using_Symbolic_States}. Their applicability often depends on the target property class or program form; \toolname\ uses LLMs to cover the evaluated C/ACSL patterns.

These static lines make different trade-offs. Constraint methods can use
Farkas' lemma and quantifier elimination to solve for invariant coefficients;
polyhedral abstract interpretation uses widening to ensure convergence.
Interpolation derives properties from reachability queries, while recurrence
methods seek closed forms. Polynomial invariants have also been synthesized
through semidefinite programming~\cite{Synthesizing_Invariants_for_Polynomial_Programs_by_Semidefinite_Programming}.
Such methods provide precise guarantees within their supported forms but may
require templates or manual choices that limit broader applicability.

\smallskip
\noindent\textbf{Machine learning and SMT-based approach.}
Learning-based invariant generators include CODE2INV~\cite{Code2Inv_a_Deep_Learning_Framework_for_Program_Verification}, which combines graph representations with CEGIS~\cite{Counterexample_Guided_Inductive_Synthesis_Modulo_Theories} and reinforcement learning; LIPuS~\cite{Loop_Invariant_Inference_through_SMT_Solving_Enhanced_Reinforcement_Learning}, which couples RL and SMT with a reward for sparse feedback; and CLN2INV~\cite{Ryan2020CLN2INV}, which uses continuous logic networks. Other learning approaches include~\cite{iRank,Lemur,LLM_Generated_Invariants_for_Bounded_Model_Checking_Without_Loop_Unrolling}. CLAUSE2INV~\cite{DBLP:journals/pacmse/CaoWXYWCM25} combines LLM-proposed clauses with counterexamples; Loopy~\cite{kamath2023findinginductiveloopinvariants} filters LLM-generated ACSL invariants with Houdini and error-guided repair. These tools primarily infer goal-driven loop invariants, often for numeric programs, rather than goal-independent functional contracts. Table~\ref{tab:comparison} compares the available SyGuS/NLA/OOPSLA results; the numeric tools do not report OOPSLA results.

CODE2INV's reinforcement learning faces sparse rewards; LIPuS explicitly
separates candidate generation from SMT checking and changes the reward
structure to address that problem. CLAUSE2INV instead generates clauses and
combines them using counterexamples, and Loopy moves this style of search to
ACSL annotations checked by Frama-C. Despite their different search
strategies, a supplied verification goal guides the invariant they seek.
This matters for our setting, where the tool must choose functional properties
to specify even when no assertion is present.

\smallskip
\noindent\textbf{LLM-based approach.}
Pei \emph{et al.}~\cite{Can_Large_Language_Models_Reason_About_Program_Invariants} fine-tune LLMs to predict invariants without formal guarantees. \specgen~\cite{SpecGen} generates Java/JML candidates and repairs failures through mutation. Preguss~\cite{Wang2025ATO} uses abstract-interpretation RTE alarms to synthesize minimal safety specifications. For C/ACSL, \autospec~\cite{autospec} decomposes programs and prompts with examples, while ACInv~\cite{Enhancing_Automated_Loop_Invariant_Generation_for_Complex_Programs_with_Large_Language_Models} feeds loop and data-structure abstractions to the LLM. Section~\ref{sec:experiments} compares \toolname\ with \autospec\ and ACInv.

Their validation goals differ. Pei \emph{et al.} assess predicted invariants
without a proof obligation; \specgen\ seeks verifiable contracts but operates
in Java/JML. Preguss focuses on preconditions sufficient for runtime-error
freedom, whereas \autospec\ and ACInv target C/ACSL contracts using verifier
feedback. These distinctions determine which results can be compared in the
same verifier setting and which serve only as cross-language context.

\begin{table*}[t]
\centering
\small
\renewcommand{\arraystretch}{0.95}
\setlength{\tabcolsep}{4pt}
\caption{Positioning of \toolname\ against representative specification and invariant generation tools.}
\label{tab:positioning}
\begin{adjustbox}{max width=\textwidth}
\begin{tabular}{lllllll}
\toprule
\textbf{Tool} & \textbf{Target} & \textbf{Goal} & \textbf{Method} & \textbf{Program Support} & \textbf{Array/Struct/Rec.} & \textbf{Validation} \\
\midrule
Daikon & invariants & test-driven & traces/templates & Java/C/C++ traces & support/no opt. & trace filter \\
LIPuS & loop inv. & goal-driven & SMT/RL clauses & C loops & unsupported & SMT \\
CLN2INV & loop inv. & goal-driven & continuous logic nets & C loops & unsupported & SMT \\
CLAUSE2INV & loop inv. & goal-driven & LLM clauses + CEX & C loops & unsupported & SMT \\
Loopy & loop inv. & goal-driven & LLM + Houdini repair & C loops & unsupported & Frama-C \\
\autospec & ACSL specs & goal-driven & decomposition + few-shot & C function/loops & support/no opt. & Frama-C \\
ACInv & ACSL specs & goal-driven & static abstractions & C function/loops & optimized & Frama-C \\
Preguss & ACSL specs & goal-driven & RTE alarms & C function/loops & support/no opt. & Frama-C \\
\specgen & JML specs & \rev{goal-independent} & LLM + mutation & Java function/loops & support/no opt. & OpenJML \\
\toolname & ACSL specs & \rev{goal-independent} & symbolic states/templates & C function/loops & optimized & Frama-C \\
\bottomrule
\end{tabular}
\end{adjustbox}
\end{table*}

ACInv and \toolname\ both supply analysis facts to LLMs. ACInv uses loop and data-structure abstractions; \toolname's symbolic snapshots additionally preserve path conditions, memory updates, alias distinctions, and field correlations needed for path-sensitive contracts. Preguss instead targets RTE freedom. \toolname\ generates functional summaries and uses RTE reports as warnings or, when execution is blocked, as a reason to fall back; refinement may strengthen safety preconditions.

The distinction is especially relevant for branches and pointer updates:
an abstract description of a loop or data structure may identify its shape
without retaining the condition under which a particular post-state relation
holds. Symbolic paths give the generator those conditional relations, while
the fallback and verifier-guided refinement handle cases outside the exact
symbolic fragment. This is a difference in the information supplied to the
model, not a claim that symbolic execution removes all proof limitations.

\smallskip
\noindent\textbf{Lean-based specification generation.}
CLEVER~\cite{thakur2025clever} and VERINA~\cite{ye2026verina} generate specifications and proofs for programs written in Lean. Lean supports expressive inductive predicates and proofs, but applying these systems to existing C code requires rewriting the implementation. \toolname\ instead attaches ACSL contracts to native C.

\smallskip
\noindent Across these lines, goal-driven invariant tools target narrower properties; \autospec\ and ACInv use C/ACSL but depend on verification goals; \specgen\ is goal-independent in Java/JML; and Lean systems require a Lean implementation. \toolname\ targets goal-independent, verifier-checked functional specifications for native C (Table~\ref{tab:positioning}).
